\chapter{Considerações finais}
\label{sec:conclusao}

A proibição da pulverização aérea no Ceará é uma política de saúde com quase
seis anos de vigência plena, um vizinho que não a adotou e nenhuma avaliação. Este
projeto propõe avaliá-la com um desenho de fronteira na Chapada do Apodi,
organizado pela teoria da mudança da lei: os elos observáveis da cadeia, do voo
ao nascimento, têm cada um a sua medida, e o efeito estimado só é atribuído à
hipótese da política depois de examinados os caminhos concorrentes.

O efeito esperado sobre a média municipal é pequeno, e o desenho principal não
o detecta. A avaliação pode entregar um limite para o efeito médio, a medida de
quanto avião a lei removeu, a fração da população ao alcance da deriva e o
comportamento dos demais elos: substituição de tecnologia, intoxicações,
produção e perda fetal.

Da avaliação do desenho sai uma recomendação que não depende dos resultados: a
política carece de monitoramento. Os registros de aplicação e de receituário
permitiriam acompanhar o cumprimento da lei e a substituição de tecnologia, e a
exceção para drones, de 2024, é ocasião para instituí-lo.

\section{Limitações}

\begin{itemize}
\item A exposição da gestante não é observada. Todos os efeitos são de
intenção de tratar no nível do município.
\item Há de dois a três municípios tratados no núcleo. A inferência depende de
métodos para poucos grupos tratados, que têm baixo poder.
\item O primeiro estágio depende de pedidos de acesso à informação, que podem
ser negados ou devolver registros incompletos.
\item A linha de base do método de aplicação é de 2006.
\item Um choque estadual simultâneo à lei só é fechado pelas especificações
secundárias.
\item A análise não é inteiramente cega. A série estadual de produção e os
números de dois municípios tratados são públicos e foram vistos. Uma versão
anterior deste projeto, com outro desenho e só com municípios cearenses,
estimou efeitos sobre o peso ao nascer depois de 2019.
\end{itemize}
